\chapter{Project Structure}
\section{Repository Organization}
KAVACH separates its REST backend from its browser frontend. The following tree shows the application directories relevant to the implementation. Generated build output, dependency folders, and local configuration values are omitted.
\begingroup\footnotesize\setstretch{1.0}
\begin{lstlisting}[basicstyle=\ttfamily\scriptsize]
UserLoginProject/
  login-analysis-backend/
    pom.xml, mvnw, mvnw.cmd
    src/main/java/com/loginanalysis/
      LoginAnalysisBackendApplication.java
      config/       JWT filter, security, passwords, AI setup
      controller/   REST entry points
      dto/          Request and response contracts
      entity/       JPA model
      repository/   Queries and aggregates
      service/      Business logic
      agent/
        TwoLoopAgentService.java
        model/      Roles, messages, tools, execution steps
        tools/      SecurityToolsRegistry.java
      util/         Provider and request helpers
      exception/    Application errors
    src/main/resources/application.properties
    src/test/java/com/loginanalysis/
      LoginAnalysisBackendApplicationTests.java
  login-analysis-frontend/
    package.json, vite.config.js
    public/
    src/
      App.jsx, main.jsx
      api/          HTTP client and feature APIs
      context/      Authentication and theme
      pages/        Feature pages
      components/   Layout and reusable UI
      utils/        Logging and animation
      assets/
\end{lstlisting}
\endgroup

\clearpage\section{Principal Implementation Modules}
Table~\ref{tab:appendix-modules} lists implementation responsibilities. The backend package root is shown below; frontend paths start at \code{login-analysis-frontend/src/}.
\par{\footnotesize\code{login-analysis-backend/src/main/java/com/loginanalysis/}}\par
\begingroup\footnotesize\setstretch{1.0}\setlength{\parskip}{0pt}\renewcommand{\arraystretch}{1.1}
\begin{longtable}{>{\raggedright\arraybackslash}p{3.8cm}>{\raggedright\arraybackslash}p{10.3cm}}
\caption{Principal application modules and responsibilities.}\label{tab:appendix-modules}\\\toprule Responsibility & Modules\\\midrule\endfirsthead
\toprule Responsibility & Modules\\\midrule\endhead
Authentication & \code{service/AuthService.java}, \code{service/JwtService.java}, \code{config/PasswordConfig.java}.\\
Access control & \code{config/SecurityConfig.java}, \code{config/JwtAuthenticationFilter.java}, \code{service/UserAuthorizationService.java}.\\
Telemetry and detection & \code{service/LoginLogService.java}, \code{service/AnomalyDetectionService.java}, \code{repository/LoginLogRepository.java}.\\
Risk and analytics & \code{service/RiskService.java}, \code{service/LoginAnalyticsService.java}, \code{service/SecurityDashboardService.java}, \code{repository/AnomalyRepository.java}.\\
User and audit operations & \code{service/UserService.java}, \code{service/AuditLogService.java}, and their controllers/repositories.\\
AI explanations & \code{service/GrokAIService.java}, \code{util/GrokClient.java}, \code{config/GrokConfig.java}.\\
Agent execution & \code{agent/TwoLoopAgentService.java}, \code{agent/tools/SecurityToolsRegistry.java}, \code{agent/model/AgentTool.java}.\\
API contracts & \code{controller/AIAnalysisController.java} and the request/response classes in \code{dto/}.\\
Frontend state and routing & \code{App.jsx}, \code{context/AuthContext.jsx}, \code{context/ThemeContext.jsx}, \code{api/client.js}.\\
Frontend views & \code{pages/LoginPage.jsx}, \code{DashboardPage.jsx}, \code{AnomaliesPage.jsx}, \code{RiskPage.jsx}, \code{LoginLogsPage.jsx}, \code{UsersPage.jsx}, \code{AuditPage.jsx}, and \code{GrokAiPage.jsx} within \code{pages/}.\\\bottomrule
\end{longtable}\endgroup

The backend follows a controller--service--repository organization. Request DTOs define inputs, services coordinate business operations, and repositories query mapped entities. The frontend uses a shared Axios client and common authentication state. The agent registry is an additional service entry path; it must be considered alongside direct controller calls when assessing authorization.
